\medskip
\noindent\emph{Preparation with the RV coordinates retained.}
By the standing convention, write
$B\subseteq\VF^{d_0}\times\RV_\infty^{q_0}$.
First replace $A,A'$ by the graphs of $f,f'$.
Their projections back to $A,A'$ are definable bijections,
and the graphs still belong to $\VF_*$.
On each compactness piece, write
$\bm U_b=(U_b,p_b)$ and $\bm U'_b=(U'_b,p'_b)$,
with fixed label degrees $k,k'$.
All the RV coordinates of $u\in U_b$, including the finite
tags, remain part of the representative.

Apply the preceding preparation lemma to the finite list of
VF polynomials appearing in the defining formulas for $B$,
the representative sets and their structure-map graphs.
These polynomials involve the VF base variables and ground
coefficients; the RV parameters remain free variables.
Thus preparation can be performed on the VF projection of
the base. Carry along the original RV coordinates and every
new RV coordinate introduced by the preparation.
RV contractibility of the prepared polynomials makes all
these RV-sorted data depend on the prepared VF base only
through its RV tuple.

After reparametrizing, and partitioning the zero-coordinate
strata, write the base as
\[
 B^\flat=\{(c,t):\tau=(\rv(c),t)\in S\},\qquad
 S\subseteq(\RV^\times)^d\times\RV_\infty^q.
\]
Here $c$ has $d$ nonzero VF coordinates and $t$ contains
all the retained RV coordinates; the finite stratum tags
are retained as well.
The inverse preparation recovers the original parameter
$b=b(c,t)$.
The representatives and label maps are now denoted by
\[
 \bm U_\tau=(U_\tau,p_\tau),\qquad
 p_\tau:U_\tau\longrightarrow(\RV^\times)^k,
\]
and similarly for the primed family.
Using the transported lift bijections, the first total space
is definably bijective to
\[
 Y=\{(c,t,x,u):
       \tau=(\rv(c),t)\in S,\ u\in U_\tau,\
       \rv(x)=p_\tau(u)\}.
\]
The primed construction is identical.
These total spaces belong to $\VF_*$: they are obtained by
the displayed bijections from the original graph families,
and RV-fiber dimension is invariant under definable
bijections.

Put $\tau=(s,t)$ and form the RV object
\[
 W=\{(\tau,u):\tau\in S,\ u\in U_\tau\},\qquad
 q(\tau,u)=(s,p_\tau(u)).
\]
Its structure map $q$ is finite-to-one.
Indeed, fix $(s,w)$ and choose any
$c\in\rv^{-1}(s)$ and $x\in\rv^{-1}(w)$.
Then $q^{-1}(s,w)$ identifies, by keeping $(t,u)$, with the
RV fiber of $Y$ at the fixed VF tuple $(c,x)$.
That fiber is finite because $Y\in\VF_*$.
The choices of $c,x$ are used only to prove finiteness and
are not parameters in the definition of $W$ or $q$.
Thus $\bm W=(W,q)\in\RV[d+k]$.

There is an explicit definable bijection
\[
 \mathbb L\bm W\longrightarrow Y,\qquad
 ((c,x),((s,t),u))\longmapsto(c,t,x,u).
\]
Its inverse restores the forced coordinate $s=\rv(c)$.
Define $\bm W'$ in the same way, of degree $d+k'$.
Consequently the two global integrals are represented by
$\bm W,\bm W'$ modulo the lifting congruence.
The coordinates $t$ remain RV coordinates; they are not
replaced by VF coordinates or given additional lifts.

This preparation does not descend the fiber RES equality
from $\mathbb S\langle b(c,t)\rangle$ to
$\mathbb S\langle\tau\rangle$.
That equality is still interpreted over the former base.
Its descent requires the finite RES-relation witnesses in
the later argument; once those witnesses have been
uniformized, their formulas can be included in the same
finite preparation package.

\medskip
\noindent\emph{Controlling the value parameters.}
For the subsequent RV calculations set $U=W$, $U'=W'$,
\[
 D=\vrv(U),\qquad D'=\vrv(U'),\qquad E=\vrv(S).
\]
Work componentwise on fixed zero-coordinate strata,
retaining their finite tags.
For $\tau\in S$, let $D_\tau=\vrv(U_\tau)$,
where the value map is taken on all remaining nonzero
coordinates of the underlying representative.
The label maps are retained separately.

We claim that a finite ground-definable partition
$S=\coprod_a S_a$ makes $D_\tau$ depend only on
$\gamma=\vrv(\tau)$ on each $S_a$.
More precisely, there are ground-definable DOAG sets
$\mathcal D_a\subseteq\vrv(S_a)\times\Gamma^r$ with
\[
 D_\tau=(\mathcal D_a)_{\vrv(\tau)}
       \qquad(\tau\in S_a).
\]
The same partition can be chosen for the primed family.

Here are the parameter details.
Stable embeddedness of $\Gamma$, followed by compactness,
first gives finitely many fixed DOAG formulas
$\Phi_\ell(\delta,z)$ such that every $D_\tau$ has the form
$\Phi_\ell(\Gamma^r,c)$ for some tuple $c$ in $\Gamma$.
Partition $S$ according to the first applicable formula.
For a fixed formula the set
\[
 P_\tau=\{c:\ \forall\delta,\
       (\delta\in D_\tau\ \Longleftrightarrow\
                    \Phi_\ell(\delta,c))\}
\]
is a nonempty definable set in a power of $\Gamma$.

Fix a ground-definable $\epsilon>0$ in $\Gamma$.
There is a definable selection $c(\tau)\in P_\tau$ using
only $\tau$ and $\epsilon$.
For completeness, in one ordered coordinate choose $0$
when available. Otherwise use the first positive component,
if one exists, reflecting the negative part if necessary.
Choose its left endpoint if included, the midpoint for an
open bounded interval, or its left endpoint plus $\epsilon$
for an open ray. The component endpoints are definable
from the set. Apply this rule successively to projections
and fibers for tuples. O-minimality makes these choices
uniformly definable. Hence the selection does not retain
any externally chosen Gamma parameter.

For a nonzero RV tuple $\tau$, RV quantifier elimination
gives the parameter-control identity
\[
 \dcl(\mathbb S,\tau)\cap\Gamma
 =
 \operatorname{span}_{\Q}
 \bigl(\dcl(\mathbb S)\cap\Gamma,\ \vrv(\tau)\bigr).
\]
To see the point, on each leading-term guard the value
conditions are rational linear conditions in the value
variables, with constants from the ground and the values
of RV monomials in $\tau$.
A uniquely defined Gamma element is a boundary value of
these conditions, and therefore has the displayed rational
affine form.
Apply this to every coordinate of $c(\tau)$.

These rational-affine descriptions cover every ambient
$\tau$; compactness makes the list finite.
Refine $S$ by their equality loci and by the earlier formula
choices. On each resulting piece,
\[
 c(\tau)=L_a\vrv(\tau)+\alpha_a,
 \qquad L_a\text{ a rational matrix},\quad
 \alpha_a\in(\dcl(\mathbb S)\cap\Gamma)^{\dim c}.
\]
Consequently one may take
\[
 \mathcal D_a=
 \{(\gamma,\delta):\gamma\in\vrv(S_a),\
             \Phi_{\ell(a)}(\delta,L_a\gamma+\alpha_a)\}.
\]
This is defined over the original ground.
Intersect the two finite partitions to handle the primed
family simultaneously.
On a selected base piece the notation
$D_{\vrv(\tau)}$, $D'_{\vrv(\tau)}$ is therefore justified.
The first partition may depend on angular data in $\tau$;
only after that partition do the value fibers depend on
$\vrv(\tau)$ alone.
